
\chapter{Análise do Código MATLAB}
\author{Expert Software Engineer}
\section{Visão Geral}

\subsection{Objetivo Principal}

The primary purpose of this MATLAB file is to experiment with \textbf{neural networks applied to algebraic/combinatorial structures represented in the complex plane}. The program explores mappings between discrete symbolic structures and complex numbers, and attempts to \textbf{learn the inverse mapping} using a neural network trained via backpropagation.

The currently active configuration corresponds to:

\[
\texttt{edited = '25Apr2025'}
\]

which executes the most recent experiment in the code.

\subsection{Tipo de Análise Performed}

The file performs the following computational tasks:

\begin{itemize}
\item Construction of a discrete combinatorial space \(A\) consisting of sequences of symbols.
\item Mapping these sequences into the complex plane using a structured embedding.
\item Training a neural network to recover the original symbolic sequence from its complex representation.
\item Evaluating prediction accuracy and geometric reconstruction error.
\end{itemize}

The model is trained using:

\begin{itemize}
\item GELU activation functions
\item Softmax output layer
\item Cross-entropy style gradients
\item Gradient descent weight updates
\end{itemize}

\subsection{Mathematical / Theoretical Context}

The mathematical framework combines several areas:

\begin{itemize}
\item \textbf{Discrete combinatorics}: sequences \(A = B^X\)
\item \textbf{Complex embeddings}: mapping symbolic sequences into complex numbers
\item \textbf{Representation theory / algebraic encodings}: using exponentials and powers to encode symbolic information
\item \textbf{Machine learning}: neural networks attempting to invert the embedding
\item \textbf{Geometric representation}: analyzing the structure of embeddings in the complex plane
\end{itemize}

Conceptually, the program studies whether a neural network can \textbf{decode algebraic information embedded in complex-valued geometric representations}.

\newpage

\section{Activation Funções}

\subsection{Explanation}

This section defines the activation functions used in the neural network.

\textbf{GELU (Gaussian Erro Linear Unit)} is used as the hidden-layer activation function, which is common in modern neural architectures.

The derivative is explicitly implemented to support backpropagation.

A numerically stable \textbf{softmax} function is defined for the output layer.

\textbf{Entradas:}
\begin{itemize}
\item Vetor \(x\) or \(z\)
\end{itemize}

\textbf{Saídas:}
\begin{itemize}
\item Activated values
\end{itemize}

\subsection{Code Snippet}

\begin{lstlisting}[language=Matlab]
gelu = @(x) 0.5 .* x .* (1 + tanh(sqrt(2/pi) .* (x + 0.044715 .* x.^3)));

gelu_derivative = @(x) ...
    0.5 .* (1 + tanh(sqrt(2/pi)*(x + 0.044715.*x.^3))) + ...
    0.5 .* x .* (1 - tanh(sqrt(2/pi)*(x + 0.044715.*x.^3)).^2) .* ...
    sqrt(2/pi) .* (1 + 3*0.044715*x.^2);

softmax = @(z) exp(z - max(z)) ./ sum(exp(z - max(z)));
\end{lstlisting}

\newpage

\section{Forward Aprovado Função}

\subsection{Explanation}

The \textbf{forward} function performs inference through the neural network.

The network structure is:

\[
\text{Entrada} \rightarrow W_1 \rightarrow \text{Hidden Layer} \rightarrow W_2^{(1..n)} \rightarrow \text{Hidden Stack} \rightarrow W_3 \rightarrow \text{Softmax}
\]

Key details:

\begin{itemize}
\item First hidden layer uses GELU activation.
\item A sequence of \(n_{lay}\) hidden layers follow.
\item Final output layer produces probabilities via softmax.
\end{itemize}

\textbf{Entradas}

\begin{itemize}
\item \(x\): embedded input vector
\item \(W_1,W_2,W_3\): weight matrices
\item \(b_1,b_2,b_3\): biases
\end{itemize}

\textbf{Saídas}

\begin{itemize}
\item \(y_3\): predicted class probabilities
\item intermediate activations for backpropagation
\end{itemize}

\subsection{Code Snippet}

\begin{lstlisting}[language=Matlab]
function [y3,y2,y1,z3,z2,z1]=forward(x,W1,W2,W3,b1,b2,b3)

nlay=size(W2,3);
nlin=size(W2,1);

y2=zeros(nlin,nlay);
z2=y2;

z1 = W1 * x + b1;
y1 = gelu(z1);

z2(:,1) = W2(:,:,1) * y1 + b2(:,1);
y2(:,1) = gelu(z2(:,1));

for i=2:nlay
    z2(:,i) = W2(:,:,i) * y2(:,i-1) + b2(:,i);
    y2(:,i) = gelu(z2(:,i));
end

z3 = W3 * y2(:,end) + b3;
y3 = softmax(z3);

end
\end{lstlisting}

\newpage

\section{Backpropagation Função}

\subsection{Explanation}

The \textbf{backprop} function computes gradients and updates network parameters using gradient descent.

The algorithm:

\begin{enumerate}
\item Compute output error using softmax-crossentropy derivative
\item Propagate gradients through hidden layers
\item Compute weight gradients
\item Update weights using learning rate \(\eta\)
\end{enumerate}

\textbf{Saídas}

Updated weight matrices and biases.

\subsection{Code Snippet}

\begin{lstlisting}[language=Matlab]
dz3 = y3 - y_true;

dW3 = dz3 * y2(:, end)';
db3 = dz3;

dy2 = W3' * dz3;

for i = nlay:-1:2
    dz2(:, i) = dy2 .* gelu_derivative(z2(:, i));
    input_to_this_layer = y2(:, i - 1);

    dW2(:, :, i) = dz2(:, i) * input_to_this_layer';
    db2(:, i) = dz2(:, i);

    dy2 = W2(:, :, i)' * dz2(:, i);
end
\end{lstlisting}

\newpage

\section{Desempenho Evaluation}

\subsection{Explanation}

This function evaluates prediction quality after each training epoch.

Metrics computed:

\begin{itemize}
\item Number of correctly predicted symbols
\item Mean geometric reconstruction error
\item Maximum geometric reconstruction error
\end{itemize}

The geometric error compares the predicted embedding against the true complex embedding.

\subsection{Code Snippet}

\begin{lstlisting}[language=Matlab]
correct_symbols = sum(a_list == a0_list, 2);
total_correct = sum(correct_symbols);

x0_pred = sum(geom(a0_list), 2);

geom_error = abs(x0_list - x0_pred);

fprintf('Epoch %d Results:\n', epoch);
fprintf('Total correct symbols: %d\n', total_correct);
\end{lstlisting}

\newpage

\section{Combinatorial Space Construction}

\subsection{Explanation}

This section constructs the symbolic space \(A\).

\[
A = (B+B+B)^X
\]

where

\begin{itemize}
\item \(nX = 3\) sequence length
\item \(nB = 4\) base alphabet
\end{itemize}

Two bijections are defined:

\begin{itemize}
\item index $\rightarrow$ sequence
\item sequence $\rightarrow$ index
\end{itemize}

\subsection{Code Snippet}

\begin{lstlisting}[language=Matlab]
nX = 3; 
nB = 4; 
nA = (3*nB)^nX;

i2s = @(i) mod(floor(bsxfun(@rdivide, i(:)-1, nB.^((1:nX)-1))), nB) + 1;
s2i = @(s) (s-1)*(nB.^((1:nX)'-1)) + 1;
\end{lstlisting}

\newpage

\section{Complex Embedding}

\subsection{Explanation}

Each symbolic sequence \(a\) is mapped to a complex value.

The embedding uses:

\[
\rho(a) = 3^{\text{position weight}}
\]

and

\[
\exp\left(\frac{2\pi i a}{n_B}\right)
\]

Thus the geometric representation becomes:

\[
\sum \rho(a) e^{2\pi i a/n_B}
\]

This generates a unique geometric encoding of each sequence.

\subsection{Code Snippet}

\begin{lstlisting}[language=Matlab]
rho=@(a) 3.^(repmat(1:nX,size(a,1),1)-1-floor(nX/2)-floor((a-1)/nB)*nX);

geom=@(a) rho(a).*exp(2i*pi*a/nB);
\end{lstlisting}

\newpage

\section{Training Loop}

\subsection{Explanation}

The neural network is trained over multiple epochs.

For each symbolic sequence:

\begin{enumerate}
\item Compute complex embedding
\item Generate feature vector
\item Run forward pass
\item Predict symbol probabilities
\item Apply backpropagation
\end{enumerate}

The training uses a \textbf{three-stage decoding process}:

\begin{itemize}
\item nano stage
\item micro stage
\item macro stage
\end{itemize}

Each stage refines the predicted symbol.

\subsection{Code Snippet}

\begin{lstlisting}[language=Matlab]
for epoch=1:Nepoch

for ix=1:nA

a=i2s(ix);

x0=sum(geom(a));

for j=1:nX

x=split(embed(x0,a0));

[y3,y2,y1,z3,z2,z1]=forward(x,W1,W2,W3,b1,b2,b3);

[W1, W2, W3, b1, b2, b3] =
backprop(x, y_true, y3, y2, y1, z3, z2, z1,
         W1, W2, W3, b1, b2, b3, eta);

end
end
end
\end{lstlisting}

\newpage

\section{Saídas e Elementos Visuais Esperados}

The program generates several graphical outputs.

\subsection{Complex Embedding Visualization}

A scatter plot in the complex plane:

\[
z = \sum \rho(a)e^{2\pi ia/n_B}
\]

\textbf{Axes}

\begin{itemize}
\item X-axis: real part
\item Y-axis: imaginary part
\end{itemize}

Each point represents one symbolic sequence.

\subsection{Training Desempenho Gráficos}

Two plots summarize training performance.

\textbf{Plot 1}

Prediction accuracy over epochs.

\[
\text{Accuracy} =
\frac{\text{correct symbols}}{\text{total symbols}}
\]

\textbf{Plot 2}

Geometric reconstruction errors:

\begin{itemize}
\item Mean error
\item Maximum error
\end{itemize}

These show how well the predicted sequence reconstructs the original complex embedding.
